%%%%%%%%%%%%%%%%%%%%%%%%%%%%%%%%%%%%%%%%%%%%%%%%%%%%%%%%%%%%%%%%%%%%%
%%%%%%%%%%%%%%%%%%%%%%%%%%%%%%%%%%%%%%%%%%%%%%%%%%%%%%%%%%%%%%%%%%%%%
%%%%%%%%%%%%%%%%%%%%%%%%%%%%%%%%%%%%%%%%%%%%%%%%%%%%%%%%%%%%%%%%%%%%%
%%%%%%%%%%%%%%%%%%%%%%%%%%%%%%%%%%%%%%%%%%%%%%%%%%%%%%%%%%%%%%%%%%%%%
\chapter{NAT/Firewall}\index{NAT!problem}
%%%%%%%%%%%%%%%%%%%%%%%%%%%%%%%%%%%%%%%%%%%%%%%%%%%%%%%%%%%%%%%%%%%%%
%%%%%%%%%%%%%%%%%%%%%%%%%%%%%%%%%%%%%%%%%%%%%%%%%%%%%%%%%%%%%%%%%%%%%
%%%%%%%%%%%%%%%%%%%%%%%%%%%%%%%%%%%%%%%%%%%%%%%%%%%%%%%%%%%%%%%%%%%%%
%%%%%%%%%%%%%%%%%%%%%%%%%%%%%%%%%%%%%%%%%%%%%%%%%%%%%%%%%%%%%%%%%%%%%

We use RFC 3022 \cite{RFC-3022-NAT} (Traditional IP Network Address Translator) as the main reference in this chapter.

Network Address Translation (NAT) is a technique for overcoming the 32-bit limit of the IPv4 address space in the TCP/IP protocol design.
It allows IP addresses to be reused within a private local network.
Communication with the outside world is achieved through a gateway with two network interfaces: one with a private IP address and the other with a public IP address.

The gateway manages the translation process, also known as NAT, which is depicted in Figure \ref{fig:NatArchitecture}.
The main element of this mechanism is \textbf{a mapping table}, which stores information about the mapping of packet flows between private addresses and ports on the two interfaces.
The NAT mechanism modifies the packet contents in the following fields:
\begin{paklist}
\item source address and port (when a packet moves from the private to the public network),
\item destination address and port (in the reverse direction, from the public to the private network),
\item checksum (recalculated over the modified packet).
\end{paklist}

When a packet leaves the private network, its \textbf{source address} is replaced with the gateway's \textbf{public address}.
The receiver can see only the public address and therefore sends its response to the gateway's IP address.
On the return path, the packet's destination IP address is translated again according to the mapping table and the packet is delivered to the original node inside the private network.
The NAT mechanism must recalculate the packet checksum because the packet has been modified.

\begin{figure}
\centering
\includesvg[width=1\textwidth]{nat-architecture}
\caption{NAT architecture.}
\label{fig:NatArchitecture}
\end{figure}

%%%%%%%%%%%%%%%%%%%%%%%%%%%%%%%%%%%%%%%%%%%%%%%%%%%%%%%%%%%%%%%%%%%%%
%%%%%%%%%%%%%%%%%%%%%%%%%%%%%%%%%%%%%%%%%%%%%%%%%%%%%%%%%%%%%%%%%%%%%
\section{NAT-related issues}
%%%%%%%%%%%%%%%%%%%%%%%%%%%%%%%%%%%%%%%%%%%%%%%%%%%%%%%%%%%%%%%%%%%%%
%%%%%%%%%%%%%%%%%%%%%%%%%%%%%%%%%%%%%%%%%%%%%%%%%%%%%%%%%%%%%%%%%%%%%
Address translation modifies only the IP packet header and does not touch the payload.
It is clear that such a mechanism cannot handle protocols that include routing information in the payload, a limitation that is even more significant when packets are encrypted.

NAT allows communication to be initiated by nodes inside the private network. Communication initiated in the opposite direction is not directly possible.
Although this is a drawback, the principle can also be considered a security feature.
In our case, however, we would like to connect nodes transparently. We will describe the QuickBridge tool later in this thesis.
Packets arriving at the gateway from the public network are forwarded only if a corresponding entry exists in the NAT mapping table.
One possible workaround is to initialise the mapping entry by periodically sending dummy packets to the outside network (keep-alive).
This solution, however, does not work with all types of NAT.

%%%%%%%%%%%%%%%%%%%%%%%%%%%%%%%%%%%%%%%%%%%%%%%%%%%%%%%%%%%%%%%%%%%%%
%%%%%%%%%%%%%%%%%%%%%%%%%%%%%%%%%%%%%%%%%%%%%%%%%%%%%%%%%%%%%%%%%%%%%
\section{A simple NAT-traversal test}
%%%%%%%%%%%%%%%%%%%%%%%%%%%%%%%%%%%%%%%%%%%%%%%%%%%%%%%%%%%%%%%%%%%%%
%%%%%%%%%%%%%%%%%%%%%%%%%%%%%%%%%%%%%%%%%%%%%%%%%%%%%%%%%%%%%%%%%%%%%
The NAT traversal mechanism can be tested as follows.
Our test configuration consists of three components: a gateway, a server and a client.

\paragraph*{NAT gateway:}
A Linux machine with two network cards: \verb|$NAT_PUB_DEV|, with IP address \verb|$NAT_PUB_IP|, connected to the public network, and \verb|$NAT_PRIV_DEV|, with IP address \verb|$NAT_PRIV_IP|, connected to the private network. NAT is enabled on the \verb|$NAT_PUB_DEV| interface using the \textbf{iptables} tool.

\paragraph*{Server with a public IP address:}
The public IP address is \verb|$SERVER_IP|. For testing, we use \textbf{netcat}, which can create a UDP server listening on port \verb|$SERVER_PORT|. The contents of each packet are written to standard output (stdout), while data read from standard input (stdin) is sent as a UDP packet to the same address and port.

\paragraph*{Client with a private IP address:}
The client is located in the private network and has IP address \verb|$CLIENT_IP|. Its network interface parameters are assigned automatically using DHCP, and its \textbf{default gateway} is set to \verb|NAT_PRIV_IP|. We again use \textbf{netcat} to facilitate UDP communication through stdin and stdout.

\bigskip
The first UDP packet sent from the node with a private IP address to the server with a public IP address creates a path for subsequent UDP communication in both directions. As described above, NAT translates packets according to the mapping entry that has been created.
This path usually disappears when the NAT gateway is restarted or when the mapping entry expires due to a timeout.

\begin{verbatim}
# gateway: enable NAT
iptables -t nat -A POSTROUTING -j MASQUERADE -o $NAT_PUB_DEV

# server: test connectivity
ping $NAT_PUB_IP

# server: listen for UDP packets
netcat -lup $SERVER_PORT

# client: connect to the server
netcat -u $SERVER_IP $SERVER_PORT
\end{verbatim}

%%%%%%%%%%%%%%%%%%%%%%%%%%%%%%%%%%%%%%%%%%%%%%%%%%%%%%%%%%%%%%%%%%%%%
%%%%%%%%%%%%%%%%%%%%%%%%%%%%%%%%%%%%%%%%%%%%%%%%%%%%%%%%%%%%%%%%%%%%%
\section{Tunneling}
%%%%%%%%%%%%%%%%%%%%%%%%%%%%%%%%%%%%%%%%%%%%%%%%%%%%%%%%%%%%%%%%%%%%%
%%%%%%%%%%%%%%%%%%%%%%%%%%%%%%%%%%%%%%%%%%%%%%%%%%%%%%%%%%%%%%%%%%%%%

In the previous test, we demonstrated a simple NAT traversal solution in which communication is initiated from within the private network and then continues bidirectionally over a single TCP or UDP port.
If we could encapsulate our communication in a single packet stream, we would be able to \textbf{tunnel} arbitrary communication between two nodes.
Such a hypothetical tunnel could provide the following services:
\begin{paklist}
\item encrypted communication,
\item reliable delivery,
\item transparent tunneling,
\item tunneling of multicast packets,
\item tunneling of other protocols, such as IPX,
\item and more \ldots
\end{paklist}

For multimedia streaming, an unreliable delivery service is usually the best solution. Additional layers, such as encryption, could have a negative impact on system latency.
Through AGConnector, AccessGrid supports multicast tunneling using the UMTP protocol.

\subsection{NAT-to-NAT communication}

So far, we have considered situations in which a client in a NAT-protected network communicates with a public server.
However, it is also possible to address the case of two NAT-protected nodes located in separate networks.

In the following example, the \textbf{nat-traverse} tool is used together with the PPP protocol to create a tunnel. NAT on both sides must preserve port numbers.
The example shows a NAT gateway \verb|\${CLIENT_IP}| for the client and a NAT gateway \verb|\${CLIENT_IP}| for the server.
Both parties exchange UDP packets through port 40000.

\begin{verbatim}
# client: create the ppp0 interface
nat-traverse --cmd="pppd updetach noauth passive notty ipparam \
  vpn 10.0.0.1:10.0.0.2" 40000:\${CLIENT_IP}:40000
\end{verbatim}

\begin{verbatim}
# server: accept the PPP connection
nat-traverse --cmd="pppd nodetach notty noauth" \
  40000:\${SERVER_IP}:40000
\end{verbatim}

After the tunnel has been established, routing tables must be configured on both sides.
QuickBridge packets can be sent through the tunnel. Alternatively, multicast packets can be tunnelled, but this requires the installation of a multicast router.
The configuration also requires administrative privileges so that the PPP network interface (typically \verb|ppp0|) can be created.

\begin{verbatim}
# client: route multicast packets through the tunnel
route add -net 224.0.0.0 netmask 255.0.0.0 gw ppp0
route add -net 239.0.0.0 netmask 255.0.0.0 gw ppp0
\end{verbatim}

\begin{verbatim}
# server: enable packet forwarding and start the multicast router
echo 1 > /proc/sys/net/ipv4/ip_forward
mrouted
\end{verbatim}

The nat-traverse application can be used with OpenVPN as well as PPP, allowing the entire communication to be encrypted.
\proofed
